\documentclass[11pt,a4paper]{article}
\usepackage[utf8]{inputenc}
\usepackage[T1]{fontenc}
\usepackage[margin=2.3cm]{geometry}
\usepackage{booktabs}
\usepackage{longtable}
\usepackage{array}
\usepackage{fancyhdr}
\usepackage{xcolor}
\usepackage[hidelinks]{hyperref}
\usepackage{parskip}

\definecolor{uteqverde}{RGB}{27,94,60}
\definecolor{ambar}{RGB}{176,122,0}
\definecolor{rojo}{RGB}{150,30,30}

\pagestyle{fancy}
\fancyhf{}
\fancyhead[L]{\footnotesize UTEQ · ISR-401 · CarniCore}
\fancyhead[R]{\footnotesize Autoevaluación FAIR · versión 2.2}
\fancyfoot[C]{\footnotesize Página \thepage}
\renewcommand{\headrulewidth}{0.4pt}

\newcommand{\si}{\textcolor{uteqverde}{\textbf{Cumple}}}
\newcommand{\parcial}{\textcolor{ambar}{\textbf{Parcial}}}
\newcommand{\no}{\textcolor{rojo}{\textbf{No cumple}}}

\begin{document}

\begin{center}
{\Large\bfseries\color{uteqverde} Autoevaluación FAIR del paquete de datos}\\[4pt]
{\large Sistema CarniCore · Entrega 4 (2B) · versión corregida 2.2}\\[8pt]
{\small Modelo aplicado: FAIR Data Maturity Model, Research Data Alliance (2020)}\\
{\small Objeto evaluado: \url{https://doi.org/10.5281/zenodo.22225854}}\\
{\small Fecha de evaluación: 4 de septiembre de 2026}
\end{center}

\vspace{4pt}
\hrule
\vspace{10pt}

\section*{1. Por qué esta versión corrige a la anterior}

La autoevaluación adjunta a la entrega del 1 de septiembre de 2026 se adjudicaba
\textbf{16 de 16 indicadores (100\,\%)}. Esa puntuación no era defendible por
tres razones concretas, señaladas en el informe docente:

\begin{enumerate}
\item El manifiesto de integridad \texttt{checksums.sha256} fallaba en 7 de sus
      86 entradas. Un paquete cuya integridad no verifica no puede declararse
      perfectamente reutilizable.
\item El informe describía el contenido del paquete de forma inexacta: declaraba
      90 respuestas de cuestionario cuando el archivo real contiene 31, un
      resultado del detector de 1/27 cuando la salida real es 0/27, y un
      orquestador \texttt{run\_all.sh} que no existe en el repositorio.
\item El informe anterior daba por accesible el pre-registro citando el
      identificador del proyecto en lugar del identificador del registro. El
      registro es en efecto público (\texttt{yp7t3}), pero el documento no
      permitía verificarlo.
\end{enumerate}

Esta versión evalúa los mismos 16 indicadores con criterio conservador. El
resultado es \textbf{14 de 16 (87,50\,\%)}, por encima del mínimo del 60\,\% que
exige la Sección 7.5 de la guía, y con los dos indicadores no cumplidos
identificados y con plan de subsanación.

\section*{2. Resultado por dimensión}

\begin{center}
\begin{tabular}{lccc}
\toprule
\textbf{Dimensión} & \textbf{Cumplidos} & \textbf{Total} & \textbf{Porcentaje} \\
\midrule
Findable (localizable)     & 5 & 5 & 100,0\,\% \\
Accessible (accesible)     & 4 & 4 & 100,0\,\% \\
Interoperable              & 2 & 3 & 66,7\,\% \\
Reusable (reutilizable)    & 3 & 4 & 75,0\,\% \\
\midrule
\textbf{Total}             & \textbf{14} & \textbf{16} & \textbf{87,50\,\%} \\
\bottomrule
\end{tabular}
\end{center}

\section*{3. Evaluación indicador por indicador}

\begin{longtable}{@{}p{1.1cm}p{5.3cm}p{2.1cm}p{7.3cm}@{}}
\toprule
\textbf{ID} & \textbf{Indicador} & \textbf{Veredicto} & \textbf{Evidencia o motivo} \\
\midrule
\endhead
F1 & Identificador global único y persistente & \si & DOI 10.5281/zenodo.22225854 \\
F2 & Metadatos descriptivos ricos & \si & Título, autoría con ORCID, palabras clave, descripción y licencia en el registro Zenodo \\
F3 & Los metadatos incluyen el identificador del dato & \si & El DOI figura en \texttt{CITATION.cff}, en \texttt{README.md} y en el manuscrito \\
F4 & Registro en un recurso indexable & \si & Zenodo indexa en OpenAIRE y DataCite \\
F5 & Identificadores de las personas autoras & \si & ORCID de los cinco integrantes \\
\midrule
A1 & Recuperable por protocolo estándar y abierto & \si & HTTPS sobre Zenodo, sin autenticación \\
A2 & Protocolo abierto, gratuito y universal & \si & Descarga directa, sin registro previo \\
A3 & Metadatos accesibles aunque el dato deje de estarlo & \si & Zenodo conserva el registro de metadatos de forma persistente \\
A4 & Metadatos del protocolo accesibles y consultables & \si & Registro público en OSF, \url{https://osf.io/yp7t3}, con estado \textit{Public registration} y marca de tiempo del 2 de agosto de 2026 \\
\midrule
I1 & Representación formal y accesible del conocimiento & \si & JSON, CSV y Markdown; esquema documentado en \texttt{README\_dataset.md} \\
I2 & \textbf{Vocabularios controlados con identificadores} & \parcial & Las palabras clave siguen el vocabulario de la revista objetivo, pero el esquema del corpus de RF no se mapea a ningún vocabulario formal (por ejemplo, SKOS o un perfil de aplicación). Se contabiliza como no cumplido \\
I3 & Referencias cruzadas cualificadas a otros datos & \si & Enlaces recíprocos entre Zenodo, OSF, GitHub y Software Heritage \\
\midrule
R1 & Atributos descriptivos plurales y relevantes & \si & Diccionario de datos completo por archivo \\
R2 & Licencia de uso clara y accesible & \si & CC BY 4.0 para datos, MIT para código, con alcance explícito en \texttt{LICENSE} \\
R3 & Procedencia detallada del dato & \si & \texttt{ANONYMIZATION.md}, \texttt{ETHICS.md}, protocolo, desviaciones y pipeline versionado \\
R4 & \textbf{Cumple estándares de comunidad del dominio} & \parcial & El paquete no adopta un estándar de metadatos del dominio de la ingeniería de software empírica (por ejemplo, el perfil de artefactos de ACM/SIGSOFT ni las guías de \textit{empirical standards}). Se contabiliza como no cumplido \\
\bottomrule
\end{longtable}

\section*{4. Indicadores no cumplidos y plan de subsanación}

\paragraph{A4 --- Resuelto.}
El registro se identificó y se verificó como público: \url{https://osf.io/yp7t3},
estado \textit{Public registration}, marca de tiempo del 2 de agosto de 2026.
La consulta previa fallaba porque se estaba interrogando el identificador del
\textit{proyecto} (\texttt{wud69}) y no el del \textit{registro}
(\texttt{yp7t3}), que son objetos distintos en OSF. Todos los documentos del
repositorio citan ya el identificador correcto.

\paragraph{I2 --- Vocabulario controlado del corpus.}
\textit{Responsable:} Pérez Ruiz. \textit{Plazo:} versión siguiente del paquete.
Documentar el esquema de \texttt{rf27.json} y de
\texttt{clasificaciones\_detector.csv} como un perfil de aplicación explícito,
declarando tipos, cardinalidades y valores admitidos de cada campo, y
publicarlo junto al dataset.

\paragraph{R4 --- Estándares del dominio.}
\textit{Responsable:} Quintero Gende. \textit{Plazo:} versión siguiente.
Adecuar el paquete a los criterios de artefacto evaluable de la comunidad de
ingeniería de software empírica (estructura, \texttt{INSTALL}, \texttt{STATUS} y
guion de reproducción de un solo comando). El pipeline ya se ejecuta con
\texttt{python run\_all.py}, de modo que el trabajo restante es
fundamentalmente documental.

\section*{5. Verificación de integridad asociada}

A diferencia de la versión anterior, esta evaluación se emite después de haber
comprobado la integridad del árbol:

\begin{center}
\begin{tabular}{lr}
\toprule
\textbf{Comprobación} & \textbf{Resultado} \\
\midrule
Entradas en \texttt{checksums.sha256} & 599 \\
Entradas que verifican & 599 (100\,\%) \\
Hashes de contenedores frente a \texttt{fichas\_tecnicas.csv} & \textbf{18 de 18 coinciden} \\
\bottomrule
\end{tabular}
\end{center}

\vspace{6pt}
\noindent\textbf{Nota de corrección (4 de septiembre de 2026).} Una
verificación previa había señalado 3 discrepancias (\texttt{06}, \texttt{07}
y \texttt{15-evidencias\_restringidas.7z}). Dos de ellas (\texttt{06},
\texttt{07}) eran un falso positivo: el hash coincidía, solo estaba escrito
en mayúsculas en \texttt{fichas\_tecnicas.csv} frente a minúsculas en el
cálculo de comparación. La tercera (\texttt{15}) era una discrepancia real:
el equipo verificó con \texttt{Get-FileHash} sobre el archivo del
repositorio que el contenido corresponde a la entrevista de ENTR-15 y que el
hash y el tamaño declarados habían quedado desactualizados frente al
contenedor actual (probablemente por una recompresión o resubida posterior
al cálculo original). Se actualizó la fila correspondiente en
\texttt{fichas\_tecnicas.csv} con el hash y tamaño reales
(\texttt{039b8def...177b1}, 24\,158\,760 bytes) y quedó verificado.

Reproducible con \texttt{bash herramientas/regenerar\_checksums.sh --check}.

\vspace{14pt}
\noindent\rule{\textwidth}{0.4pt}\\
{\footnotesize Documento generado desde \texttt{fair\_assessment.tex} y
versionado en la raíz del repositorio. Autoevaluación conforme al modelo de la
Research Data Alliance (2020); no sustituye a una ejecución de la herramienta
F-UJI contra el DOI, que se recomienda adjuntar cuando el registro OSF esté
publicado.}

\end{document}
